% ============================================================
% sections/sensitivity_analysis.tex — 灵敏度分析 / Sensitivity Analysis
% ============================================================
% Test how results change when key parameters vary.
% This is a CORE requirement for math modeling papers.
% 这是数学建模论文的核心要求。

\section{灵敏度分析}

\subsection{分析目的}

[说明进行灵敏度分析的目的：检验模型的稳定性和鲁棒性]

\subsection{分析方法}

[说明采用的分析方法：固定其他参数，逐个改变关键参数，观察结果的变化趋势。]

\subsection{参数灵敏度分析}

\subsubsection{参数 $\alpha$ 的灵敏度分析}

固定其他参数不变，将 $\alpha$ 在 $[\alpha_{\min}, \alpha_{\max}]$ 范围内变化，观察目标函数值的变化。

\begin{table}[htbp]
\centering
\caption{参数 $\alpha$ 的灵敏度分析结果}
\label{tab:sens_alpha}
\begin{tabular}{cc}
\toprule
\textbf{$\alpha$ 取值} & \textbf{目标函数值} \\
\midrule
[$\alpha_1$] & [目标值1] \\
[$\alpha_2$] & [目标值2] \\
[$\alpha_3$] & [目标值3] \\
\bottomrule
\end{tabular}
\end{table}

% Figure placeholder example / 图片占位示例
% \begin{figure}[htbp]
%     \centering
%     \includegraphics[width=0.7\textwidth]{figures/sensitivity_alpha.png}
%     \caption{参数 $\alpha$ 的灵敏度曲线}
%     \label{fig:sens_alpha_curve}
% \end{figure}

\subsubsection{参数 $\beta$ 的灵敏度分析}

[对 $\beta$ 进行类似分析]

\subsection{分析结论}

[总结灵敏度分析的主要结论：模型对哪些参数敏感，哪些参数不敏感，模型的适用范围]
